% ============================================================================
% Supplementary Information SI-6: Shared Geometry and Optimisation Machinery
%
% Drop-in subsection fragment. Parent SI must provide: amsmath/amssymb,
% booktabs, siunitx (with \angstrom), and the tcolorbox environments
% {evaluation}, {concern}, {threshbox}. Bibliography handled at the
% parent level. (Parameter reference lives in si_parameters.tex.)
% ============================================================================

\subsection{Shared Machinery: ML Force Field, Geometry, and Optimisation}
\label{si:shared}

The discovery loops (SI-\ref{si:pes}, SI-\ref{si:generative}) and the kinetics pipeline (SI-\ref{si:kinetics}) share a single on-line energy oracle and a small set of numerical primitives that are not specific to any one loop. This subsection consolidates them. Chemistry-specific thresholds (e.g., the asymmetric \SI{0.04}{}/\SI{0.005}{eV} IRC endpoint criterion in SI-\ref{si:pes}, or the \SI{1.0}{\angstrom} RMSD endpoint match in SI-\ref{si:generative}) stay in the calling SI; this section defines the model and the algorithms those thresholds gate.

%-----------------------------------------------------------------------------
\subsubsection{ML force field (MD-ET)}
\label{sec:shared-mdet}
%-----------------------------------------------------------------------------

All energies, forces, and Hessians used by the on-line pipeline come from a single transformer-based neural network potential, \textbf{MD-ET}. It is a \emph{direct force model}: forces are predicted as a per-atom 3-vector output head rather than obtained by differentiating an energy. Energy and force are independent heads, and the Hessian is obtained by automatic differentiation of the force head with respect to atomic positions (one row per force component, batched via \texttt{torch.vmap}). Radial cutoff for the neighbour list: $\SI{5.0}{\angstrom}$.

Three depths are released as the \texttt{variant} parameter --- \texttt{4l}, \texttt{5l}, \texttt{12l} --- with \texttt{12l} the production default. The model is trained on the QCML dataset (broad small-organic chemical space, DFT-level reference energies and forces); the transition-state subset is not represented separately, which is the structural reason the per-saddle Hessian still benefits from the trust-region machinery below and from the validation gates in SI-\ref{si:pes}/SI-\ref{si:generative}. The TS-proposer diffusion model (SI-\ref{si:generative} Stage 3) is trained on a separate TS-relaxed dataset and is independent of the MD-ET training distribution.

Distribution and loading: the model ships as the \texttt{md-et} pip package and is loaded via \texttt{md\_et.load\_calculator(variant=...)}; weights are pulled from the Hugging Face Hub on first use and the worker Dockerfile pre-warms the cache at build time. The returned object exposes the standard ASE \texttt{Calculator} interface plus a bespoke \texttt{get\_batched\_hessians} for the throughput-critical Hessian path. The default Hessian batch size is $16$; the production deployment uses $8$ to stay within L4 24~GB VRAM at typical pipeline sizes (the calculator's built-in halving recovery only applies within a single call, not across calls in a batch). A \texttt{filter\_forces=true} option enables a noise-suppression filter on the predicted forces and is on across the production deployment.

Routing every on-line ML call through a single direct-force calculator means the ML in-box barriers consumed by the kinetics solver (SI-\ref{si:kinetics}), P-RFO convergence, IRC endpoint stability, and the generative critic's barrier gates all depend on the same model. The DFT energy refinement (\S\ref{sec:kinetics-dft}) provides the independent ground truth that the kinetics ultimately consumes, so MD-ET acts as a fast filter whose role is to keep the candidate set tractable; the final rate constants depend on DFT whenever the CPU pipeline has caught up.


%-----------------------------------------------------------------------------
\subsubsection{Permutation-invariant RMSD}
\label{sec:shared-rmsd}
%-----------------------------------------------------------------------------

Comparing two molecular geometries requires aligning them in space and resolving the labelling ambiguity for chemically equivalent atoms. The system uses Kabsch alignment for rigid-body rotation and a two-strategy resolution for atom permutations:

\begin{enumerate}
    \item \textbf{Kabsch alignment.} Centre both structures at the origin and compute the optimal rotation via singular-value decomposition of the cross-covariance:
    \begin{equation}
        \mathbf{R}^* = \arg\min_{\mathbf{R} \in SO(3)} \sum_i \lVert \mathbf{r}_i^{(1)} - \mathbf{R}\,\mathbf{r}_i^{(2)} \rVert^2.
    \end{equation}
    The SVD-based form handles the right-handed-frame correction automatically (no reflections introduced).

    \item \textbf{Permutation resolution.} Define $\sigma$ over atoms of identical atomic number. RMSD is computed as
    \begin{equation}
        \text{RMSD} = \min_{\sigma}\ \sqrt{\frac{1}{N}\,\sum_i \lVert \mathbf{r}_i^{(1)} - \mathbf{R}^* \mathbf{r}_{\sigma(i)}^{(2)} \rVert^2}.
    \end{equation}
    The minimisation strategy depends on problem size:
    \begin{itemize}
        \item If the total number of permutations is $\leq 720 = 6!$, brute-force enumerate all of them and re-run Kabsch for each. This is needed when one element dominates the structure (e.g.\ $\mathrm{HCO_3^-}$ has 3 of 5 atoms as oxygen): the initial Kabsch alignment is then unreliable because most atoms are interchangeable, and Hungarian assignment on a single bad alignment gets stuck in a local minimum.
        \item Otherwise: Hungarian linear-assignment algorithm on the distance matrix from the Kabsch-aligned configuration. For larger molecules ($> 6$ equivalent atoms per element) the alignment from non-equivalent atoms is good enough that Hungarian finds the optimum on the first pass.
    \end{itemize}
\end{enumerate}

\paragraph{Canonical atom ordering.}
Before any comparison the atoms are reordered into the RDKit canonical-rank order. This makes the RMSD deterministic across workers and across different input orderings.

\paragraph{Where used.}
SI-\ref{si:pes} Stage 6 (endpoint distinctness $> \SI{0.25}{\angstrom}$); SI-\ref{si:generative} Stage 6 (endpoint-to-source match $< \SI{1.0}{\angstrom}$); minimum deduplication in the per-compound PES graph (RMSD tiers).


%-----------------------------------------------------------------------------
\subsubsection{Connected-component fragment detection}
\label{sec:shared-fragments}
%-----------------------------------------------------------------------------

A molecular geometry is decomposed into chemically-meaningful fragments by building an adjacency matrix from a covalent-radius-scaled distance criterion,
\begin{equation}
    \text{adj}_{ij} = \mathbb{I}\bigl[\lVert r_i - r_j \rVert < 1.3\,(r^{\text{cov}}_i + r^{\text{cov}}_j)\bigr],
\end{equation}
with covalent radii from ASE. Connected components are extracted by depth-first search and returned as sorted index tensors. The $1.3\times$ scale is conservative; values around $1.15$--$1.20\times$ are more typical for cheminformatics but would mis-classify the stretched bonds at un-converged IRC endpoints as already-dissociated.

Per-fragment validity rules used by both loops: a fragment is accepted iff it is non-empty, its composition satisfies the tetrose cap $\{C{:}5,\ H{:}10,\ O{:}5\}$, and (when single-atom) the atom is hydrogen.

\paragraph{Where used.}
SI-\ref{si:generative} Stage 7 (product-side fragmentation detection); SI-\ref{si:generative} Stage 6 greedy validator (fragmenting both endpoints for chemical-difference comparison); SI-\ref{si:pes} Stage 7 (escape-vs-intramolecular classification by SMILES of relaxed endpoints).


%-----------------------------------------------------------------------------
\subsubsection{SMILES canonicalisation and charge inference}
\label{sec:shared-smiles}
%-----------------------------------------------------------------------------

A 3D geometry is converted to a canonical SMILES via RDKit's bond perception (\texttt{rdkit.Chem.rdDetermineBonds.DetermineBonds}) followed by \texttt{Chem.MolToSmiles(\dots, isomericSmiles=False)}. The non-isomeric form is used everywhere so stereoisomers collapse to the same string --- the diffusion proposer cannot reliably distinguish stereoisomers in any case, and the graph would otherwise admit spurious duplicates that differ only in inverted chirality of a single carbon.

Integer charge is inferred from the perceived bond pattern by summing formal charges atom-by-atom after RDKit's valence-completion pass. The charge is used at three points: (i)~when registering a discovered compound, (ii)~when checking charge conservation across a reaction, and (iii)~when handing a structure to the ML calculator. The DFT refinement (SI-\ref{si:kinetics}) re-derives charges from a wavefunction calculation downstream.


%-----------------------------------------------------------------------------
\subsubsection{Hessian computation and projection}
\label{sec:shared-hessian}
%-----------------------------------------------------------------------------

The Hessian is the central object for both saddle search and validation. The full $3N \times 3N$ matrix is computed by automatic differentiation of the ML forces with respect to atomic positions (\S\ref{sec:shared-mdet}, \texttt{get\_batched\_hessians}), and is then projected to remove the six trivial modes (three translations, three rotations) that contaminate the spectrum at non-zero geometric centring or non-zero overall angular momentum:
\begin{equation}
    \mathbf{H}_\text{proj} = \mathbf{P}\,\mathbf{H}\,\mathbf{P}, \qquad \mathbf{P} = \mathbf{I} - \sum_{k=1}^{6} \mathbf{v}_k \mathbf{v}_k^\top,
\end{equation}
where $\{\mathbf{v}_k\}$ is a mass-weighted orthonormal basis of the translation/rotation subspace obtained by Gram--Schmidt on the six analytical T/R vectors. The projected Hessian is diagonalised into $\{\lambda_i, \mathbf{u}_i\}$ pairs (sorted ascending); eigenvalues with $|\lambda_i| < 10^{-4}\,\text{eV/\AA}^2$ are masked as numerical-noise floor before any chemical interpretation. The two loops apply \emph{different} TS-spectrum policies and we do not unify them: the PES loop (\S\ref{sec:pes-prfo}) admits only \emph{exactly one} significant negative eigenvalue, while the generative loop (\S\ref{sec:hessian-irc}) admits \emph{at least one} below $\SI{-0.01}{eV/\angstrom^2}$. This is a deliberate precision/throughput trade-off: the PES loop trusts its P-RFO walker to converge to a clean first-order saddle, whereas the generative loop accepts higher-order Hessian signatures from diffusion-proposed geometries because the IRC and barrier gates downstream filter out the consequences. Curvature ($\lambda < \SI{-0.01}{eV/\angstrom^2}$) and significance ($|\lambda| < 10^{-4}\,\text{eV/\AA}^2$) thresholds also differ by loop and are stated in the calling SIs.


%-----------------------------------------------------------------------------
\subsubsection{FIRE optimiser}
\label{sec:shared-fire}
%-----------------------------------------------------------------------------

The Fast Inertial Relaxation Engine \citep[][used force-only]{bitzek2006fire} is a velocity-Verlet-based first-order minimiser. Each step:
\begin{enumerate}
    \item Velocity-Verlet update of positions and velocities under the current forces.
    \item Project velocity onto force direction: $\mathbf{v} \leftarrow (1-\alpha)\mathbf{v} + \alpha\,|\mathbf{v}|\,\hat{\mathbf{F}}$.
    \item If $\mathbf{F} \cdot \mathbf{v} > 0$ (descent direction) for the past $N_\text{min}$ steps, increase $\Delta t$ by a factor $f_\text{inc}$ (up to $\Delta t_\text{max}$) and decay $\alpha$ by $f_\alpha$; otherwise reset $\mathbf{v} = 0$, reset $\alpha$ to its initial value, and shrink $\Delta t$ by $f_\text{dec}$.
\end{enumerate}
Default constants follow the original paper: $\Delta t_\text{init} = 0.1$, $\Delta t_\text{max} = 1.0$, $\alpha_\text{start} = 0.1$, $N_\text{min} = 5$. FIRE is the workhorse for any optimisation that does not need a Hessian: per-fragment relax (SI-\ref{si:generative} Stage 7), merge encounter prep (SI-\ref{si:generative} Stage 2), and the NEB backbone (\S\ref{sec:shared-neb}).


%-----------------------------------------------------------------------------
\subsubsection{Trust-region Newton minimiser}
\label{sec:shared-newton}
%-----------------------------------------------------------------------------

When a Hessian is available --- e.g., during IRC endpoint relaxation (SI-\ref{si:pes} Stage 6, SI-\ref{si:generative} Stage 5) --- a second-order trust-region Newton step is used instead of FIRE:
\begin{enumerate}
    \item Project out T/R modes (\S\ref{sec:shared-hessian}), diagonalise, and floor the absolute eigenvalues to $|\lambda_\text{min}| = 10^{-2}\,\text{eV/\AA}^2$ to prevent pathologically large steps in flat (floppy) directions.
    \item Compute the rational-function downhill step in each mode, $p_i = g_i / \lambda_i$, where $g_i$ is the gradient projection on the $i$-th eigenvector.
    \item Constrain the total step to the current trust radius (clamp on per-mode magnitudes if the global norm exceeds the radius).
    \item Evaluate the trial position, compute $\rho = \Delta E_\text{actual} / \Delta E_\text{predicted}$ where $\Delta E_\text{predicted} = -\mathbf{f}\cdot\mathbf{p} + \tfrac{1}{2}\mathbf{p}^\top\mathbf{H}\mathbf{p}$, and accept if $\rho > 0.1$. On rejection, shrink the trust radius by a factor of $2$ and retry (up to 5 retries before giving up).
    \item After acceptance: shrink the trust radius if $\rho < 0.25$ or $\rho > 2.0$ (poor quadratic model), expand by $1.5$ if $0.5 < \rho < 1.5$ and the step hit the trust boundary (good model).
\end{enumerate}
Default trust radius for endpoint relaxation: $\Delta_0 = \SI{0.3}{\angstrom}$, bounded by $[\SI{0.005}{\angstrom}, \SI{2.0}{\angstrom}]$. P-RFO (SI-\ref{si:pes} Stage 4) uses a \emph{partitioned} trust region with a separate TS-mode radius scaled by $1.5\times$ the minimisation radius; that machinery is loop-specific and stays in SI-\ref{si:pes}.


%-----------------------------------------------------------------------------
\subsubsection{Intrinsic Reaction Coordinate (IRC)}
\label{sec:shared-irc}
%-----------------------------------------------------------------------------

Both loops validate a candidate transition state by following the imaginary normal mode in both directions and relaxing each endpoint to a minimum. The shared algorithm:

\begin{enumerate}
    \item Project the Hessian (\S\ref{sec:shared-hessian}); confirm that at least one significant eigenvalue is negative (below the loop-specific curvature threshold; \S\ref{sec:pes-saddle-validation}, \S\ref{sec:hessian-irc}), and take $\hat{\mathbf{u}}_\text{TS}$ to be the eigenvector of the most negative eigenvalue. The IRC step itself does \emph{not} count negative modes beyond this. The strict ``exactly one'' criterion that defines a first-order saddle is enforced \emph{upstream} of the IRC in the PES loop only --- at P-RFO convergence (\S\ref{sec:pes-prfo}), before any candidate reaches this stage --- so by construction the geometries that enter the PES-loop IRC carry a single negative mode. The generative loop has no such upstream guarantee: its candidates come from a diffusion proposer (\S\ref{sec:diffusion}) and may carry $\geq 1$ negative modes at the IRC entry; higher-order saddles are then filtered only implicitly, by failing the endpoint or bidirectional-barrier gates downstream (\S\ref{sec:endpoints}, \S\ref{sec:bwd-barrier}).

    \item Compute a curvature-adaptive displacement magnitude $\Delta x$ from the harmonic-approximation form
    \begin{equation}
        \Delta x = \mathrm{clip}\!\left(\frac{c}{\sqrt{|\lambda_\text{TS}|}},\ \Delta x_\text{min},\ \Delta x_\text{max}\right),
    \end{equation}
    with $c$, $\Delta x_\text{min}$, $\Delta x_\text{max}$ chosen per loop (SI-\ref{si:pes}: $c = \sqrt{0.16}$, range $[0.02, 0.5]\,\si{\angstrom}$; SI-\ref{si:generative}: $c = 0.2$, range $[0.05, 0.5]\,\si{\angstrom}$). Stiff saddles get smaller displacements; flat saddles get larger ones. Curvatures below $10^{-6}\,\text{eV/\AA}^2$ are treated as flat and use the maximum.

    \item Build $\mathbf{R}^\pm = \mathbf{R}_\text{TS} \pm \Delta x\,\hat{\mathbf{u}}_\text{TS}$ and relax each independently to a minimum using trust-region Newton (\S\ref{sec:shared-newton}). Loop-specific convergence: SI-\ref{si:pes} uses $f_\text{max} = \SI{0.005}{eV/\angstrom}$, max $120$ steps; SI-\ref{si:generative} uses the same $f_\text{max}$ but $350$ steps, reflecting the wider distribution of saddle curvatures the diffuser produces. Bail-late acceptance at $f_\text{max} < 10\times$ tolerance is shared.

    \item \emph{Bail-early rule:} if the forward direction fails to relax, skip the backward direction entirely. The saddle is rejected immediately; this saves one expensive relaxation per failure.

    \item Assign direction by energy: the higher-energy endpoint is labelled \emph{reactant} and the lower-energy \emph{product}. This stores every reaction downhill in the ML view, which collapses direction ambiguity for per-pair deduplication. SI-\ref{si:pes} additionally enforces endpoint distinctness (RMSD $> \SI{0.25}{\angstrom}$) and an asymmetric energy-drop criterion; SI-\ref{si:generative} additionally enforces endpoint identity against the source compound (strict + greedy modes).
\end{enumerate}


%-----------------------------------------------------------------------------
\subsubsection{Nudged Elastic Band (NEB)}
\label{sec:shared-neb}
%-----------------------------------------------------------------------------

NEB \citep{henkelman2000climbing} is used \emph{only} for the same-basin same-or-different decision in minimum deduplication: given two relaxed minima with $E_1, E_2$, is there a barrier between them above a threshold?

\paragraph{Band construction.}
The two endpoints are optimally aligned (Kabsch + permutation; \S\ref{sec:shared-rmsd}). The number of interior images is adaptive on the endpoint RMSD,
\begin{equation}
    N_\text{images} = \mathrm{clip}\!\left(\Bigl\lceil \tfrac{\text{RMSD}}{0.15\,\si{\angstrom}} \Bigr\rceil,\ 5,\ 12\right),
\end{equation}
and images are interpolated using the Image-Dependent Pair Potential (IDPP) method via ASE, which produces a physically reasonable initial band even for substantial rearrangements (better than linear interpolation, which routinely creates atom overlaps).

\paragraph{Two-phase optimisation.}
Total FIRE budget $N_\text{FIRE} = 150$ steps with a spring constant $k_\text{spring} = \SI{0.1}{eV/\angstrom^2}$:
\begin{itemize}
    \item \textbf{Phase 1 (regular NEB):} $2/3$ of the budget. All images relaxed onto the MEP with the standard tangent projection of the spring forces.
    \item \textbf{Early bail:} if the peak energy after Phase 1 exceeds $\SI{1.0}{eV}$ above the higher endpoint, the band is considered unphysical (bad interpolation in a system where IDPP could not find a sensible path) and the same-basin check returns ``different basins''.
    \item \textbf{Phase 2 (CI-NEB):} remaining $1/3$ of the budget with the climbing-image variant: the highest-energy image is pushed against the spring forces toward the true saddle, refining the barrier estimate without re-relaxing the rest of the band.
\end{itemize}

\paragraph{Same-basin criterion.}
\begin{equation}
    \Delta E^\ddagger = \max_i(E_i) - \max(E_1, E_2),
\end{equation}
the barrier above the \emph{higher} endpoint. If $\Delta E^\ddagger < \SI{0.05}{eV}$, the two minima are merged. The threshold corresponds to roughly $2\,k_B T$ at room temperature: any barrier smaller than this is freely crossed thermally and treating the two minima as distinct would over-resolve the conformational landscape.


%-----------------------------------------------------------------------------
\subsubsection{Trajectory force integration}
\label{sec:shared-traj}
%-----------------------------------------------------------------------------

When a relaxation trajectory is available, the barrier from the trajectory endpoint to the TS is computed by integrating the force component along the path,
\begin{equation}
    \Delta E^\ddagger \;=\; \int_{\mathbf{R}_\text{min}}^{\mathbf{R}_\text{TS}} \mathbf{F}(\mathbf{r}) \cdot d\mathbf{r}
    \;\approx\; -\,\tfrac{1}{2}\sum_{n} \bigl(\mathbf{F}_n + \mathbf{F}_{n+1}\bigr) \cdot \bigl(\mathbf{R}_{n+1} - \mathbf{R}_{n}\bigr),
\end{equation}
discretised by the trapezoidal rule over the stored trajectory frames (positions and forces, optionally with Hessians for a second-order correction). This is preferred over the endpoint energy difference $E_\text{TS} - E_\text{min}$ because the force integral accumulates over the whole path: spurious local fluctuations in the ML energy (which scale as the per-frame noise of the energy head) average out, while the force head's smoother behaviour dominates the result. When the trajectory is too short for a meaningful integral (or unavailable), the fallback is the endpoint difference.


%-----------------------------------------------------------------------------
\subsubsection{Citation notes}
\label{sec:shared-cites}
%-----------------------------------------------------------------------------

The algorithms above are standard in computational chemistry; this section consolidates them for cross-reference rather than as novel contributions. The primary references implied throughout: Kabsch (1976) for orthogonal Procrustes; Bitzek et al.\ (2006) for FIRE; Banerjee, Adams, Simons \& Shepard (1985) and Baker (1986) for P-RFO (the source-of-truth for the trust-region partition in SI-\ref{si:pes}); Henkelman \& J\'{o}nsson (2000) for the improved-tangent NEB and CI-NEB; Smidstrup et al.\ (2014) for the IDPP interpolation. Loop-specific deviations from these standard implementations (e.g., the $10$-retry overlap guard on encounter-complex displacement in SI-\ref{si:generative}, or the partitioned trust radii for P-RFO) are documented in the respective loop SI.
